% Prose rendering of PrimeTensor/Scale.lean.
% Fragment: no \documentclass, no preamble, no document environment.

\section{Intrinsic dyadic multiplicative scale}
\label{Scale:sec}

\leanfile{PrimeTensor/Scale.lean}

\subsection*{What this file establishes}

\texttt{PrimeTensor/Order.lean} left closeness parameterised by an arbitrary
tolerance $\delta$, admissible when $\delta > 1$.  This file replaces that
parameter by a single fixed ruler --- the magnitude $2$ --- and recovers
arbitrarily fine tolerances not by shrinking the ruler but by \emph{squaring
the ratio being measured}.  Closeness at level $\ell$ says that the oriented
ratio, squared $\ell - 1$ times, is still below $2$.

The substitution is exact and it is forced.  Refining a dyadic tolerance
directly would mean comparing against
$2^{1/2}, 2^{1/4}, 2^{1/8}, \dots$, and none of those is a positive rational,
so none of them is an element of $\name{MulRat}$: the tolerances one wants are
not expressible in the object language.  The equivalent condition on the
squared ratio is, and it uses only multiplication and the order already built.
This is what the module header means by admitting no rational square root.

\subsection{The ruler}

\begin{definition}[\lean{twoRatio} and \lean{MulRat.two}]\label{Scale:def:two}
\begin{align*}
  \name{twoRatio} &\defeq
    \bigl\langle\, \name{ofBag}(\ctor{factor}\;\name{primeTwo}\;\ctor{one}),
      \;1 \,\bigr\rangle
    &&: \name{PrimeRatio}, \\
  \name{two} &\defeq \name{ofRatio}(\name{twoRatio})
    &&: \name{MulRat}.
\end{align*}
\end{definition}

\begin{leancode}
def twoRatio : PrimeRatio :=
  ⟨PrimeMultiset.ofBag (.factor primeTwo .one), 1⟩

/-- Distinguished multiplicative ruler. -/
def two : MulRat :=
  ofRatio twoRatio
\end{leancode}

The numerator is the one-atom barcode carrying the prime $2$ and nothing else,
and the denominator is the pivot, so the pair denotes $2/1$.  Note that the
ruler is built, not postulated: it is an ordinary element of the carrier,
assembled from $\name{primeTwo}$ of \texttt{PrimeTensor/Prime.lean}.

\begin{lemma}[\lean{MulRat.one\_lt\_two}]\label{Scale:lem:one-lt-two}
$1 < \name{two}$ in $\name{MulRat}$.
\end{lemma}

\begin{proof}
Both sides are classes of literals, so the lifted order reduces definitionally
to the representative-level comparison
$\name{PrimeRatio.Lt}(1, \name{twoRatio})$, which by
\texttt{PrimeTensor/Ratio.lean} unfolds to
\[
  \name{eval}(1.\name{upper}) \cdot \name{eval}(\name{twoRatio}.\name{lower})
    \;<\;
  \name{eval}(\name{twoRatio}.\name{upper}) \cdot
    \name{eval}(1.\name{lower}).
\]
Simplification evaluates the four components: the unit contributes $1$ on both
sides by $\name{one\_upper}$, $\name{one\_lower}$ and
$\name{eval\_one}$, while the numerator of the ruler evaluates through
$\name{eval\_ofBag}$ and $\name{eval\_factor}$ to
$\name{primeTwo}.\name{value} \cdot 1$.  The goal becomes $1 \cdot 1 < 2 \cdot 1$,
closed by $\ctor{norm\_num}$ once $\name{primeTwo}$ is unfolded.
\end{proof}

So the ruler is an admissible tolerance in the sense of
\texttt{PrimeTensor/Order.lean}, and being fixed, it is admissible once and for
all --- no hypothesis $1 < \delta$ has to be carried anywhere below.

\subsection{Repeated squaring}

\begin{definition}[\lean{MulRat.scalePow}]\label{Scale:def:scalepow}
For $r : \name{MulRat}$, define $\name{scalePow}(r, -) : \name{Depth} \to
\name{MulRat}$ by structural recursion on the level:
\[
  \name{scalePow}(r, \ctor{one}) = r,
  \qquad
  \name{scalePow}(r, \ctor{succ}\,d) = q \cdot q
    \quad\text{where } q = \name{scalePow}(r, d).
\]
\end{definition}

\begin{leancode}
def scalePow (r : MulRat) : Depth → MulRat
  | .one => r
  | .succ d =>
      let q := scalePow r d
      q * q
\end{leancode}

\begin{proposition}[The effective exponent]\label{Scale:prop:exponent}
$\name{scalePow}(r, \ell) = r^{\,2^{\size{\ell} - 1}}$.  In particular the
exponents realised, as $\ell$ ranges over the depths, are
$1, 2, 4, 8, \dots$ --- the powers of two, and only those.
\end{proposition}

\begin{proof}
By induction on $\ell$.  For $\ell = \ctor{one}$, $\size{\ell} = 1$ and the
value is $r = r^{2^0}$.  For $\ell = \ctor{succ}\,d$, the induction hypothesis
gives $q = r^{2^{\size{d}-1}}$, and $q \cdot q = r^{2 \cdot 2^{\size{d}-1}}
= r^{2^{\size{d}}} = r^{2^{\size{\ctor{succ}\,d} - 1}}$.
\end{proof}

The definition squares the running value rather than multiplying by $r$, so
the level counts \emph{doublings of the exponent}, not the exponent itself.
Level $\ell$ therefore reaches exponent $2^{\size{\ell}-1}$ in $\size{\ell}$
steps: the ladder climbs geometrically in a linear number of levels.

\begin{lemma}[\lean{MulRat.scalePow\_one}]\label{Scale:lem:scalepow-one}
$\name{scalePow}(1, \ell) = 1$ for every level $\ell$.
\end{lemma}

\begin{leancode}
@[simp] theorem scalePow_one : ∀ d : Depth, scalePow (1 : MulRat) d = 1
  | .one => rfl
  | .succ d => by
      change scalePow (1 : MulRat) d * scalePow (1 : MulRat) d = 1
      rw [scalePow_one d, one_mul]
\end{leancode}

\begin{proof}
By structural recursion on the level.  At $\ctor{one}$ the value is the
argument, so the claim is $\ctor{rfl}$.  At $\ctor{succ}\,d$ the value is
$q \cdot q$ with $q = \name{scalePow}(1, d)$; the induction hypothesis rewrites
$q$ to $1$, and $\name{one\_mul}$ from \texttt{PrimeTensor/MulRat.lean}
finishes.  Note the recursive call is made explicitly on $d$, the lemma being
stated as a $\forall$ and proved by pattern matching rather than by
$\ctor{induction}$.
\end{proof}

The unit is thus fixed by the whole ladder, which is exactly what reflexivity
of the closeness relation will need.

\subsection{Closeness at an intrinsic scale}

\begin{definition}[\lean{MulRat.ScaleWithin}]\label{Scale:def:scalewithin}
For a level $\ell$ and $a, b : \name{MulRat}$,
\[
  \name{ScaleWithin}(\ell, a, b) \defeq
    \bigl(\name{scalePow}(\name{ratio}(a,b), \ell) < \name{two}\bigr) \wedge
    \bigl(\name{scalePow}(\name{ratio}(b,a), \ell) < \name{two}\bigr).
\]
\end{definition}

\begin{leancode}
def ScaleWithin (level : Depth) (a b : MulRat) : Prop :=
  scalePow (ratio a b) level < two ∧
  scalePow (ratio b a) level < two
\end{leancode}

\begin{proposition}[What the levels measure]\label{Scale:prop:tolerance}
$\name{ScaleWithin}(\ell, a, b)$ holds precisely when both oriented ratios lie
strictly below $2^{1/2^{\size{\ell}-1}}$.  Increasing the level therefore
tightens the requirement, and the effective tolerances
\[
  2, \; \sqrt{2}, \; \sqrt[4]{2}, \; \sqrt[8]{2}, \; \dots
\]
decrease towards $1$.
\end{proposition}

\begin{proof}
By Proposition~\ref{Scale:prop:exponent} the first conjunct is
$r^{2^{\size{\ell}-1}} < 2$ where $r = \name{ratio}(a,b)$, and for $r > 0$ this
is equivalent to $r < 2^{1/2^{\size{\ell}-1}}$, the map $x \mapsto x^{n}$ being
strictly increasing on the positive reals for $n \geq 1$.  The same applies to
the second conjunct.
\end{proof}

\begin{designnote}[Why squaring the ratio and not shrinking the ruler]
\label{Scale:note:no-roots}
Proposition~\ref{Scale:prop:tolerance} makes the design constraint visible.
The tolerances the ladder realises are $2^{1/2^{k}}$, and for $k \geq 1$ every
one of them is irrational: $\sqrt{2}$ already is.  So none of them is an
element of $\name{MulRat}$, which by construction contains exactly the positive
rationals.  A definition of the form ``$\name{ratio}(a,b) < \delta_k$'' with
$\delta_k$ the $k$-th tolerance is therefore not expressible here at all ---
not merely inconvenient, but absent from the object language.

Squaring the ratio states the same condition using only what is available: a
multiplication, the fixed ruler $\name{two}$, and the order of
\texttt{PrimeTensor/Order.lean}.  This is the sense in which the scale is
\emph{intrinsic}: the refinement is carried by the depth index, a piece of the
object language, rather than by a sequence of tolerances that would have to be
imported from outside it.

The same reading explains the other two exclusions the header names.  There is
no additive midpoint because there is no addition; and there is no zeroth
scale because levels are depths, so the coarsest tolerance is the ruler itself
and the ladder has a bottom rung rather than a degenerate one.
\end{designnote}

\begin{remark}[Level one is the old notion]\label{Scale:rem:level-one}
$\name{scalePow}(r, \ctor{one})$ is $r$ by definition, so
$\name{ScaleWithin}(\ctor{one}, a, b)$ is literally
$\name{ratio}(a,b) < \name{two} \wedge \name{ratio}(b,a) < \name{two}$, which
is $\name{Within}(\name{two}, a, b)$ in the sense of
\texttt{PrimeTensor/Order.lean} --- the two unfold to the same proposition.
The new relation is thus a refinement of the old one at a particular
tolerance, not a replacement for it.  The file does not state this.
\end{remark}

\begin{lemma}[\lean{scaleWithin\_refl} and \lean{scaleWithin\_symm}]
\label{Scale:lem:refl-symm}
$\name{ScaleWithin}(\ell, a, a)$ holds for every level $\ell$ and every $a$;
and $\name{ScaleWithin}(\ell, a, b)$ implies $\name{ScaleWithin}(\ell, b, a)$.
\end{lemma}

\begin{proof}
For reflexivity, $\name{ratio}(a,a) = 1$ by $\name{ratio\_self}$ of
\texttt{PrimeTensor/Order.lean}, then
Lemma~\ref{Scale:lem:scalepow-one} collapses the whole ladder to $1$, and each
conjunct becomes Lemma~\ref{Scale:lem:one-lt-two}.  Symmetry is the exchange of
the two conjuncts.
\end{proof}

\begin{remark}[Reflexivity is now unconditional]\label{Scale:rem:uncond}
Compare $\name{within\_refl}$ in \texttt{PrimeTensor/Order.lean}, which
required the hypothesis $1 < \delta$ and could not do without it.  Here there
is no hypothesis at all: the tolerance is $\name{two}$, fixed, and
Lemma~\ref{Scale:lem:one-lt-two} discharges the admissibility condition once
inside this file.  Every later statement about $\name{ScaleWithin}$ is free of
it.  This is a small illustration of the general point --- fixing the ruler
converts a carried hypothesis into a proved lemma.
\end{remark}

\begin{remark}[Monotonicity in the level is not proved]\label{Scale:rem:mono}
The levels ought to be nested: closeness at a finer level should imply
closeness at a coarser one, so that the relations form a decreasing family.
They do --- if $r^{2^{k-1}} \geq 2$ then squaring gives $r^{2^{k}} \geq 4 > 2$,
so contrapositively $\name{ScaleWithin}(\ctor{succ}\,d, a, b)$ implies
$\name{ScaleWithin}(d, a, b)$ --- but the file states no such lemma, and the
argument as given needs a comparison fact about $\name{MulRat}$ beyond the
order lemmas available so far.  Nor is there any composition rule across two
levels, the analogue of a triangle inequality, which
\texttt{PrimeTensor/Order.lean} also lacked.  Both are what a Cauchy notion
built on this scale will require.
\end{remark}

\subsection*{Summary of the correspondence}

\begin{center}
\small
\begin{tabular}{@{}lll@{}}
\hline
Lean declaration & Kind & Here \\
\hline
\lean{MulRat.twoRatio}          & definition & Def.~\ref{Scale:def:two} \\
\lean{MulRat.two}               & definition & Def.~\ref{Scale:def:two} \\
\lean{MulRat.one\_lt\_two}      & theorem    & Lem.~\ref{Scale:lem:one-lt-two} \\
\lean{MulRat.scalePow}          & definition & Def.~\ref{Scale:def:scalepow} \\
\lean{MulRat.scalePow\_one}     & simp thm   & Lem.~\ref{Scale:lem:scalepow-one} \\
\lean{MulRat.ScaleWithin}       & definition & Def.~\ref{Scale:def:scalewithin} \\
\lean{MulRat.scaleWithin\_refl} & theorem    & Lem.~\ref{Scale:lem:refl-symm} \\
\lean{MulRat.scaleWithin\_symm} & theorem    & Lem.~\ref{Scale:lem:refl-symm} \\
\hline
\end{tabular}
\end{center}

\noindent
Every declaration of \texttt{PrimeTensor/Scale.lean} appears above.  The file
contains no \lean{sorry}, no axiom, and no unproved named proposition.
Propositions~\ref{Scale:prop:exponent} and~\ref{Scale:prop:tolerance}, and
Remarks~\ref{Scale:rem:level-one} and~\ref{Scale:rem:mono}, are stated for the
reader and are not declarations of the file; the first two are the arithmetic
reading of the ladder, and they are what
Design note~\ref{Scale:note:no-roots} rests on.
